% Generated draft: requires mathematical and source review.
\documentclass[11pt]{article}
\usepackage[margin=25mm]{geometry}
\usepackage{fontspec}
\setmainfont{Latin Modern Roman}
\usepackage{amsmath,amssymb,mathtools,unicode-math}
\setmathfont{Latin Modern Math}
\usepackage{enumitem,needspace,xurl,hyperref}
\setlength{\parindent}{0pt}
\setlength{\parskip}{5pt}
\setlength{\emergencystretch}{3em}
\newcommand{\N}{\mathbb{N}}
\newcommand{\Z}{\mathbb{Z}}
\newcommand{\Q}{\mathbb{Q}}
\newcommand{\R}{\mathbb{R}}
\newcommand{\C}{\mathbb{C}}
\newcommand{\F}{\mathbb{F}}
\newcommand{\A}{\mathbb{A}}
\newcommand{\PP}{\mathbb{P}}
\newcommand{\E}{\mathbb{E}}
\newcommand{\eps}{\varepsilon}
\newcommand{\dd}{\,\mathrm{d}}
\begin{document}
\section*{UnsolvedMath ID: 20002031}
\textbf{Source problem code:} catalogue entry
\subsection{Definitions and mathematical statement}
The paired aid identifies this canonical record as a splice rather than a separate problem: its opening sentence finishes Problem 11, while the paragraph about order-three operators introduces Problem 12. The absent antecedent of ``this would provide a unique specification of \(Q\)'' is not reconstructed. The text mentions Rod Gover, Alice Chang and Jie Qing and an order-three \(P_3\) on a three-manifold occurring as a boundary or embedded hypersurface, with an associated \(Q_3\). It does not specify the domain, locality or normalization of the operator.

The aid supplies a locality diagnostic. On a manifold of dimension at least two, a scalar formally self-adjoint differential operator of odd order cannot be elliptic. In the usual scalar Hermitian symbol convention, its principal symbol is real and odd in the cotangent variable; connectedness of the unit cotangent sphere forces a zero. Thus an intrinsic scalar self-adjoint elliptic order-three operator cannot be a finite-order local differential operator. A pseudodifferential realization is a possible nonlocal alternative, not an identification forced by the source wording alone.

The aid distinguishes the local Chang--Qing boundary \(P_3\): its input is an ambient function near a hypersurface \(\Sigma^3\subset X^4\), and it depends on third normal jets. It does not factor through the boundary trace. Writing \(B_3\) for this type of operator to distinguish its domain, the diagnostic is that no intrinsic trace operator \(\overline P_3\) satisfies
\[
 B_3u=\overline P_3(u|_\Sigma)\qquad\text{for every ambient }u.
\]
The reported third-normal-jet dependence permits functions with zero trace but nonzero output. This boundary operator is not a counterexample to the intrinsic scalar lemma; its self-adjoint boundary theory is not automatically the same as an \(L^2(\Sigma)\) operator on arbitrary boundary traces.

For a separately specified intrinsic pair on a closed three-manifold, suppose
\[
 e^{3\omega}P_{3,\widehat g}f=P_{3,g}f,\qquad
 e^{3\omega}Q_{3,\widehat g}=Q_{3,g}+P_{3,g}\omega,
 \qquad\widehat g=e^{2\omega}g,
\]
with \(P_{3,g}\) formally self-adjoint and \(P_{3,g}1=0\). The aid reports that the affine law is consistent under composition of conformal changes and that
\[
 \int_\Sigma Q_{3,\widehat g}\,dv_{\widehat g}
 =\int_\Sigma Q_{3,g}\,dv_g.
\]
For a fixed \(P_3\), two \(Q_3\) assignments obeying this same law differ exactly by a scalar conformal invariant of weight \(-3\), meaning
\(I_{\widehat g}=e^{-3\omega}I_g\). Conversely adding any such invariant preserves the law. This is uniqueness modulo that ambiguity, not an absolute uniqueness statement or an existence theorem for the unspecified source pair.
\subsection{Short English statement}
This extraction is contextual material, not a complete independent problem. The supplied aid distinguishes an intrinsic nonlocal order-three operator from a local boundary operator acting on ambient jets. For a fixed intrinsic self-adjoint conformal pair it records total-integral invariance and the precise conformal-invariant ambiguity in \(Q_3\).
\subsection{Sources}
\begin{itemize}
\item Local source record: catalogue entry, catalogue entry; canonical \texttt{aim-geometry-notes.json}, \texttt{notes[368]}, canonical source item 11. Title: ``Distinguishing intrinsic and boundary order-three conformal operators.''
\item AIM workshop \emph{Conformal structure in geometry, analysis, and physics}, supplied reference \url{https://aimath.org/WWN/confstruct/confstruct.pdf}. The research summary identifies a splice between Problems 11 and 12. This draft uses the supplied record without accessing a PDF.
\item The source mentions Rod Gover, Alice Chang and Jie Qing; the aid specifically identifies the local Chang--Qing boundary operator. Full bibliographic citations, exact formulas and boundary conventions are absent. Attempt 1 supplies the symbol/trace diagnostic and conditional intrinsic-pair consequences.
\item Registry dated 14 August 2026: status ``open,'' classification ``OPEN-TRIAGE,'' statement status ``exact''; attempt context only, research-run status not a problem, rigor ``proved,'' novelty confidence low. The machine-generated, non-peer-reviewed aid requires independent verification before citation.
\end{itemize}
\end{document}
